%% 使用 njuthesis 文档类生成南京大学学位论文的示例文档
%%
%% 作者：胡海星，starfish (at) gmail (dot) com
%% 项目主页: http://haixing-hu.github.io/nju-thesis/
%%
%% 本样例文档中用到了吕琦同学的博士论文的提高和部分内容，在此对他表示感谢。
%%
\documentclass[macfonts,master]{njuthesis}
%% njuthesis 文档类的可选参数有：
%%   nobackinfo 取消封二页导师签名信息。注意，按照南大的规定，是需要签名页的。
%%   phd/master/bachelor 选择博士/硕士/学士论文

% 使用 blindtext 宏包自动生成章节文字
% 这仅仅是用于生成样例文档，正式论文中一般用不到该宏包
\usepackage[math]{blindtext}

%%%%%%%%%%%%%%%%%%%%%%%%%%%%%%%%%%%%%%%%%%%%%%%%%%%%%%%%%%%%%%%%%%%%%%%%%%%%%%%
% 设置《国家图书馆封面》的内容，仅博士论文才需要填写

% 设置论文按照《中国图书资料分类法》的分类编号
\classification{0175.2}
% 论文的密级。需按照GB/T 7156-2003标准进行设置。预定义的值包括：
% - \openlevel，表示公开级：此级别的文献可在国内外发行和交换。
% - \controllevel，表示限制级：此级别的文献内容不涉及国家秘密，但在一定时间内
%   限制其交流和使用范围。
% - \confidentiallevel，表示秘密级：此级别的文献内容涉及一般国家秘密。
% - \clasifiedlevel，表示机密级：此级别的文献内容涉及重要的国家秘密 。
% - \mostconfidentiallevel，表示绝密级：此级别的文献内容涉及最重要的国家秘密。
% 此属性可选，默认为\openlevel，即公开级。
\securitylevel{\controllevel}
% 设置论文按照《国际十进分类法UDC》的分类编号
% 该编号可在下述网址查询：http://www.udcc.org/udcsummary/php/index.php?lang=chi
\udc{004.72}
% 国家图书馆封面上的论文标题第一行，不可换行。此属性可选，默认值为通过\title设置的标题。
\nlctitlea{数据中心}
% 国家图书馆封面上的论文标题第二行，不可换行。此属性可选，默认值为空白。
\nlctitleb{网络模型研究}
% 国家图书馆封面上的论文标题第三行，不可换行。此属性可选，默认值为空白。
\nlctitlec{}
% 导师的单位名称及地址
\supervisorinfo{南京大学计算机科学与技术系~~南京市汉口路22号~~210093}
% 答辩委员会主席
\chairman{张三丰~~教授}
% 第一位评阅人
\reviewera{阳顶天~~教授}
% 第二位评阅人
\reviewerb{张无忌~~副教授}
% 第三位评阅人
\reviewerc{黄裳~~教授}
% 第四位评阅人
\reviewerd{郭靖~~研究员}

%%%%%%%%%%%%%%%%%%%%%%%%%%%%%%%%%%%%%%%%%%%%%%%%%%%%%%%%%%%%%%%%%%%%%%%%%%%%%%%
% 设置论文的中文封面

% 论文标题，不可换行
\title{一种支持微服务动态更新系统的设计与实现}
% 如果论文标题过长，可以分两行，第一行用\titlea{}定义，第二行用\titleb{}定义，将上面的\title{}注释掉
% \titlea{半轻衰变$D^+\to \omega(\phi)e^+\nu_e$的研究}
% \titleb{和弱衰变$J/\psi \to D_s^{(*)-}e^+\nu_e$的寻找}

% 论文作者姓名
\author{王东宇}
% 论文作者联系电话
\telphone{18362926576}
% 论文作者电子邮件地址
\email{wdongyu@outlook.com}
% 论文作者学生证号
\studentnum{MG1833071}
% 论文作者入学年份（年级）
\grade{2018}
% 导师姓名职称
\supervisor{曹春~~教授}
% 导师的联系电话
\supervisortelphone{18951679203}
% 论文作者的学科与专业方向
\major{计算机科学与技术}
% 论文作者的研究方向
\researchfield{软件方法学}
% 论文作者所在院系的中文名称
\department{计算机科学与技术系}
% 论文作者所在学校或机构的名称。此属性可选，默认值为``南京大学''。
\institute{南京大学}
% 论文的提交日期，需设置年、月、日。
\submitdate{2021年4月115日}
% 论文的答辩日期，需设置年、月、日。
\defenddate{2021年6月1日}
% 论文的定稿日期，需设置年、月、日。此属性可选，默认值为最后一次编译时的日期，精确到日。
%% \date{2013年5月1日}

%%%%%%%%%%%%%%%%%%%%%%%%%%%%%%%%%%%%%%%%%%%%%%%%%%%%%%%%%%%%%%%%%%%%%%%%%%%%%%%
% 设置论文的英文封面

% 论文的英文标题，不可换行
\englishtitle{Design and Implementation of a Microservices System Supporting Dynamic Update}
% 论文作者姓名的拼音
\englishauthor{Dongyu Wang}
% 导师姓名职称的英文
\englishsupervisor{Professor Chun Cao}
% 论文作者学科与专业的英文名
\englishmajor{Computer Science and Technology}
% 论文作者所在院系的英文名称
\englishdepartment{Department of Computer Science and Technology}
% 论文作者所在学校或机构的英文名称。此属性可选，默认值为``Nanjing University''。
\englishinstitute{Nanjing University}
% 论文完成日期的英文形式，它将出现在英文封面下方。需设置年、月、日。日期格式使用美国的日期
% 格式，即``Month day, year''，其中``Month''为月份的英文名全称，首字母大写；``day''为
% 该月中日期的阿拉伯数字表示；``year''为年份的四位阿拉伯数字表示。此属性可选，默认值为最后
% 一次编译时的日期。
\englishdate{April 15, 2021}

%%%%%%%%%%%%%%%%%%%%%%%%%%%%%%%%%%%%%%%%%%%%%%%%%%%%%%%%%%%%%%%%%%%%%%%%%%%%%%%
% 设置论文的中文摘要

% 设置中文摘要页面的论文标题及副标题的第一行。
% 此属性可选，其默认值为使用|\title|命令所设置的论文标题
\abstracttitlea{一种支持微服务动态更新系统的设计与实现}
% 设置中文摘要页面的论文标题及副标题的第二行。
% 此属性可选，其默认值为空白
% \abstracttitleb{}

%%%%%%%%%%%%%%%%%%%%%%%%%%%%%%%%%%%%%%%%%%%%%%%%%%%%%%%%%%%%%%%%%%%%%%%%%%%%%%%
% 设置论文的英文摘要

% 设置英文摘要页面的论文标题及副标题的第一行。
% 此属性可选，其默认值为使用|\englishtitle|命令所设置的论文标题
\englishabstracttitlea{Design and Implementation of a Microservices System}
% 设置英文摘要页面的论文标题及副标题的第二行。
% 此属性可选，其默认值为空白
\englishabstracttitleb{Supporting Dynamic Update}

%%%%%%%%%%%%%%%%%%%%%%%%%%%%%%%%%%%%%%%%%%%%%%%%%%%%%%%%%%%%%%%%%%%%%%%%%%%%%%%
\begin{document}

%%%%%%%%%%%%%%%%%%%%%%%%%%%%%%%%%%%%%%%%%%%%%%%%%%%%%%%%%%%%%%%%%%%%%%%%%%%%%%%

% 制作国家图书馆封面（博士学位论文才需要）
% \makenlctitle
% 制作中文封面
\maketitle
% 制作英文封面
\makeenglishtitle


%%%%%%%%%%%%%%%%%%%%%%%%%%%%%%%%%%%%%%%%%%%%%%%%%%%%%%%%%%%%%%%%%%%%%%%%%%%%%%%
% 开始前言部分
\frontmatter

%%%%%%%%%%%%%%%%%%%%%%%%%%%%%%%%%%%%%%%%%%%%%%%%%%%%%%%%%%%%%%%%%%%%%%%%%%%%%%%
% 论文的中文摘要
\begin{abstract}
复杂网络的研究可上溯到20世纪60年代对ER网络的研究。90年后代随着Internet
的发展，以及对人类社会、通信网络、生物网络、社交网络等各领域研究的深入，
发现了小世界网络和无尺度现象等普适现象与方法。对复杂网络的定性定量的科
学理解和分析，已成为如今网络时代科学研究的一个重点课题。

在此背景下，由于云计算时代的到来，本文针对面向云计算的数据中心网络基础
设施设计中的若干问题，进行了几方面的研究。………………
% 中文关键词。关键词之间用中文全角分号隔开，末尾无标点符号。
\keywords{小世界理论；网络模型；数据中心}
\end{abstract}

%%%%%%%%%%%%%%%%%%%%%%%%%%%%%%%%%%%%%%%%%%%%%%%%%%%%%%%%%%%%%%%%%%%%%%%%%%%%%%%
% 论文的英文摘要
\begin{englishabstract}
\blindtext
% 英文关键词。关键词之间用英文半角逗号隔开，末尾无符号。
\englishkeywords{Small World, Network Model, Data Center}
\end{englishabstract}

%%%%%%%%%%%%%%%%%%%%%%%%%%%%%%%%%%%%%%%%%%%%%%%%%%%%%%%%%%%%%%%%%%%%%%%%%%%%%%%
% 论文的前言，应放在目录之前，中英文摘要之后

%%%%%%%%%%%%%%%%%%%%%%%%%%%%%%%%%%%%%%%%%%%%%%%%%%%%%%%%%%%%%%%%%%%%%%%%%%%%%%%
% 生成论文目次
\tableofcontents

%%%%%%%%%%%%%%%%%%%%%%%%%%%%%%%%%%%%%%%%%%%%%%%%%%%%%%%%%%%%%%%%%%%%%%%%%%%%%%%
% 生成插图清单。如无需插图清单则可注释掉下述语句。
\listoffigures

%%%%%%%%%%%%%%%%%%%%%%%%%%%%%%%%%%%%%%%%%%%%%%%%%%%%%%%%%%%%%%%%%%%%%%%%%%%%%%%
% 生成附表清单。如无需附表清单则可注释掉下述语句。
\listoftables

%%%%%%%%%%%%%%%%%%%%%%%%%%%%%%%%%%%%%%%%%%%%%%%%%%%%%%%%%%%%%%%%%%%%%%%%%%%%%%%
% 开始正文部分
\mainmatter

%%%%%%%%%%%%%%%%%%%%%%%%%%%%%%%%%%%%%%%%%%%%%%%%%%%%%%%%%%%%%%%%%%%%%%%%%%%%%%%
% 学位论文的正文应以《绪论》作为第一章
\chapter{绪言}\label{chapter_introduction}

\section{研究背景}
\subsection{微服务架构}
随着软件系统规模的扩大，业务需求的迭代，功能的追加扩展，传统的单体应用(Monolithic)最终成为一个庞然大物，变得越来越复杂，逻辑耦合严重，难以理解。因此为了满足用户对于一定规模的软件的快速开发集成、提高可扩展性等要求，我们需要将其迁移到微服务架构\cite{bucchiarone2018monolithic,taibi2017processes}。微服务架构(Microservice Architecture)是一种基于一组独立部署运行的小型服务来构建应用的架构方法，服务间使用轻量级的通信协议进行通信\cite{Microservices}。与传统的单体应用相比，微服务架构可以帮助我们对应用进行有效地拆分，缩短软件开发周期，降低软件 维护成本，降低服务间耦合度的同时满足了软件程序对于快速持续集成和持续交付的需求\cite{balalaie2016microservices,taibi2018architectural}。由于前述的相关特性，微服务架构开始在工业领域得到广泛的应用，包括Amazon\footnote{}、Netflix、Spotify等

微服务架构的基本思想

\subsection{软件动态更新技术}

\section{研究现状}
\section{本文工作}
\section{本文组织}

%%%%%%%%%%%%%%%%%%%%%%%%%%%%%%%%%%%%%%%%%%%%%%%%%%%%%%%%%%%%%%%%%%%%%%%%%%%%%%%
\chapter{相关工作}\label{chapter_relatedwork}

%%%%%%%%%%%%%%%%%%%%%%%%%%%%%%%%%%%%%%%%%%%%%%%%%%%%%%%%%%%%%%%%%%%%%%%%%%%%%%%
\chapter{支持事务一致性的动态更新}\label{chapter_3}


%%%%%%%%%%%%%%%%%%%%%%%%%%%%%%%%%%%%%%%%%%%%%%%%%%%%%%%%%%%%%%%%%%%%%%%%%%%%%%%
\chapter{支持动态更新的Kubernetes CRD实现}\label{chapter_4}


%%%%%%%%%%%%%%%%%%%%%%%%%%%%%%%%%%%%%%%%%%%%%%%%%%%%%%%%%%%%%%%%%%%%%%%%%%%%%%%
\chapter{案例研究与实验评估}\label{chapter_experiments}

%%%%%%%%%%%%%%%%%%%%%%%%%%%%%%%%%%%%%%%%%%%%%%%%%%%%%%%%%%%%%%%%%%%%%%%%%%%%%%%
% 学位论文的正文应以《结论》作为最后一章
\chapter{总结与展望}\label{chapter_concludes}
\section{工作总结}
\section{研究展望}

%%%%%%%%%%%%%%%%%%%%%%%%%%%%%%%%%%%%%%%%%%%%%%%%%%%%%%%%%%%%%%%%%%%%%%%%%%%%%%%
% 致谢，应放在《结论》之后
\begin{acknowledgement}
  时光荏苒，研究生三年的学习时光已经接近尾声，在此我想对我的父母，我的老师和同学们表达由衷的谢意。

  感谢我的导师曹春教授。在我读研期间，曹老师一直关心我的科研和生活。本文从选题、成文到修改都离不开曹老师的耐心指导。曹老师对待学术严谨执着，待人和善，平时非常关心学生的生活状况和身体健康。他的谆谆教诲使我受益匪浅。

  感谢吕建教授、马晓星教授、陶先平教授、徐峰教授、许畅教授、黄宇教授、胡昊副教授、余萍副教授、徐经纬副教授、姚远副教授、张建莹老师等所有关心和帮助过我的老师。

  感谢实验室的所有同学，良好的学习氛围是我顺利完成科研项目的保障，能够和你们一起学习和工作让我感到非常荣幸。

  感谢同寝的邱圣广和张文明同学，你们的陪伴是我宝贵的精神财富。

  最后向我的家人和女朋友致以最诚挚的感谢。感谢父母对我的养育之恩和无条件的支持，使我可以全身心的投入到科研中；感谢我的女朋友对我的关心和鼓励，让我一路上都有人同行，充满喜悦。
\end{acknowledgement}

%%%%%%%%%%%%%%%%%%%%%%%%%%%%%%%%%%%%%%%%%%%%%%%%%%%%%%%%%%%%%%%%%%%%%%%%%%%%%%%
% 附录
\appendix

\chapter{博士(硕士)学位论文编写格式规定(试行)}

\section{适用范围}

本规定适用于博士学位论文编写，硕士学位论文编写应参照执行。

\section{引用标准}

GB7713科学技术报告、学位论文和学术论文的编写格式。

GB7714文后参考文献著录规则。

\section{印制要求}

论文必须用白色纸印刷，并用A4(210mm×297mm)标准大小的白纸。纸的四周应留足空白
边缘，上方和左侧应空边25mm以上，下方和右侧应空边20mm以上。除前置部分外，其它
部分双面印刷。

论文装订不要用铁钉，以便长期存档和收藏。

论文封面与封底之间的中缝（书脊）必须有论文题目、作者和学校名。

\section{编写格式}

论文由前置部分、主体部分、附录部分(必要时)、结尾部分(必要时)组成。

前置部分包括封面，题名页，声明及说明，前言，摘要(中、英文)，关键词，目次页，
插图和附表清单(必要时)，符号、标志、缩略词、首字母缩写、单位、术语、名词解释
表(必要时)。

主体部分包括绪论(作为正文第一章)、正文、结论、致谢、参考文献表。

附录部分包括必要的各种附录。

结尾部分包括索引和封底。

\section{前置部分}

\subsection{封面（博士论文国图版用）}

封面是论文的外表面，提供应有的信息，并起保护作用。

封面上应包括下列内容：
\begin{enumerate}
\item 分类号  在左上角注明分类号，便于信息交换和处理。一般应注明《中国图书资
  料分类法》的类号，同时应注明《国际十进分类法UDC》的类号；
\item 密级  在右上角注明密级；
\item “博士学位论文”用大号字标明；
\item 题名和副题名   用大号字标明；
\item 作者姓名；
\item 学科专业名称；
\item 研究方向；
\item 导师姓名，职称；
\item 日期包括论文提交日期和答辩日期；
\item 学位授予单位。
\end{enumerate}

\subsection{题名}

题名是以最恰当、最简明的词语反映论文中最重要的特定内容的逻辑组合。

题名所用每一词语必须考虑到有助于选定关键词和编写题录、索引等二次文献可以提供
检索的特定实用信息。

题名应避免使用不常见的缩略词、首字母缩写字、字符、代号和公式等。

题名一般不宜超过20字。

论文应有外文题名，外文题名一般不宜超过10个实词。

可以有副题名。

题名在整本论文中不同地方出现时，应完全相同。

\subsection{前言}

前言是作者对本论文基本特征的简介，如论文背景、主旨、目的、意义等并简述本论文
的创新性成果。

\subsection{摘要}

摘要是论文内容不加注释和评论的简单陈述。

论文应有中、英文摘要，中、英文摘要内容应相同。

摘要应具有独立性和自含性，即不阅读论文的全文，便能获得必要的信息，摘要
中有数据、有结论，是一篇完整的短文，可以独立使用，可以引用，可以用于推广。摘
要的内容应包括与论文同等量的主要信息，供读者确定有无必要阅读全文，也供文摘等
二次文献引用。摘要的重点是成果和结论。

中文摘要一般在1500字，英文摘要不宜超过1500实词。

摘要中不用图、表、化学结构式、非公知公用的符号和术语。

\subsection{关键词}

关键词是为了文献标引工作从论文中选取出来用于表示全文主题内容信息款目的单词或
术语。

每篇论文选取3－8个词作为关键词，以显著的字符另起一行，排在摘要的左下方。在英
文摘要的左下方应标注与中文对应的英文关键词。

\subsection{目次页}

目次页由论文的章、节、附录等的序号、名称和页码组成，另页排在摘要的后面。

\subsection{插图和附表清单}

论文中如图表较多，可以分别列出清单并置于目次页之后。

图的清单应有序号、图题和页码。表的清单应有序号、表题和页码。

符号、标志、缩略词、首字母缩写、计量单位、名词、术语等的注释表符号、标志、缩略词、
首字母缩写、计量单位、名词、术语等的注释说明汇集表，应置于图表清单之后。

\section{主体部分}

\subsection{格式}

主体部分由绪论开始，以结论结束。主体部分必须由另页右页开始。每一章必须另页开
始。全部论文章、节、目的格式和版面安排要划一，层次清楚。

\subsection{序号}

\begin{figure}[htbp]
  \centering
  \includegraphics[width= 0.5\textwidth]{njuname.eps}\\
  \caption{测试附录中的插图}\label{fig:appendix1}
\end{figure}

论文的章可以写成：第一章。节及节以下均用阿拉伯数字编排序号，如
1.1，1.1.1等。

论文中的图、表、附注、参考文献、公式、算式等一律用阿拉伯数字分别分章依序连续编排
序号。其标注形式应便于互相区别，一般用下例：图1.2；表2.3；附注1）；文献[4]；式
  (6.3)等。

论文一律用阿拉伯数字连续编页码。页码由首页开始，作为第1页，并为右页另页。封页、
封二、封三和封底不编入页码，应为题名页、前言、目次页等前置部分单独编排页码。页码
必须标注在每页的相同位置，便于识别。

\begin{equation}
    C_i = \frac{2E_i}{k_i(k_i-1)}
\end{equation}

附录依序用大写正体A、B、C、$\cdots$编序号，如：附录A。附录中的图、表、式、参考文
献等另行编序号，与正文分开，也一律用阿拉伯数字编码，但在数码前题以附条序码，如图
A.1；表B.2；式(B.3)；文献[A.5]等。

\subsection{绪论}

绪论（综述）：简要说明研究工作的目的、范围、相关领域的前人工作和知识空白、理
论基础和分析，研究设想、研究方法和实验设计、预期结果和意义等。一般在教科书中
有的知识，在绪论中不必赘述。

绪论的内容应包括论文研究方向相关领域的最新进展、对有关进展和问题的评价、本论
文研究的命题和技术路线等；绪论应表明博士生对研究方向相关的学科领域有系统深入
的了解，论文具有先进性和前沿性；

\begin{problem}
测试定理环境。测试定理环境。测试定理环境。测试定理环境。测试定理环境。测试定理环境。
测试定理环境。测试定理环境。测试定理环境。
\end{problem}

为了反映出作者确已掌握了坚实的基础理论和系统的专门知识，具有开阔的科学视野，对研
究方案作了充分论证，绪论应单独成章，列为第一章，绪论的篇幅应达$1\sim 2$万字，不
得少于$1$万字；绪论引用的文献应在$100$篇以上，其中外文文献不少于$60\%$；引用文献
应按正文中引用的先后排列。

\subsection{正文}

论文的正文是核心部分，占主要篇幅。正文必须实事求是，客观真切，准确完备，合乎
逻辑，层次分明，简便可读。

\begin{figure}[htbp]
  \centering
  \includegraphics[width= 0.5\textwidth]{njuname.eps}\\
  \caption{测试附录中的插图}\label{fig:appendix2}
\end{figure}

正文的每一章(除绪论外)应有小结，在小结中应明确阐明作者在本章中所做的工作，特
别是创新性成果。凡本论文要用的基础性内容或他人的成果不应单独成章，也不应作过
多的阐述，一般只引结论、使用条件等，不作推导。

\subsubsection{图}

图包括曲线图、构造图、示意图、图解、框图、流程图、记录图、布置图、地图、照片
、图版等。

图应具有“自明性”，即只看图、图题和图例，不阅读正文，就可以理解图意。

图应编排序号。每一图应有简短确切的图题，连同图号置于图下。必要时，应将图上的
符号、标记、代码，以及实验条件等，用最简练的文字，横排于图题下方，作为图例说
明。

\begin{example}
测试定理环境。测试定理环境。测试定理环境。测试定理环境。测试定理环境。测试定理环境。
测试定理环境。测试定理环境。测试定理环境。
\end{example}

曲线图的纵、横坐标必须标注“量、标准规定符号、单位”。此三者只有在不必要标明
(如无量纲等)的情况下方可省略。坐标上标注的量的符号和缩略词必须与正文一致。

照片图要求主题和主要显示部分的轮廓鲜明，便于制版。如用放大缩小的复制品，必须
清晰，反差适中。照片上应该有表示目的物尺寸的标度。

\subsubsection{表}

表的编排，一般是内容和测试项目由左至右横读，数据依序竖排。表应有自明性。

表应编排序号。

每一表应有简短确切的表题，连同标号置于表上。必要时，应将表中的符号、标记、代
码，以及需要说明事项，以最简练的文字，横排于表题下，作为表注，也可以附注于表
下。表内附注的序号宜用小号阿拉伯数字并加圆括号置于被标注对象的右上角，如：
xxx${}^{1)}$；不宜用“*”，以免与数学上共轭和物质转移的符号相混。

表的各栏均应标明“量或测试项目、标准规定符号、单位”。只有在无必要标注的情况下
方可省略。表中的缩略词和符号，必须与正文中一致。

表内同一栏的数字必须上下对齐。表内不宜用“同上”，“同左”和类似词，一律填入具体数字
或文字。表内“空白”代表未测或无此项，“－”或“\textellipsis”（因“－”可能与代表阴性
  反应相混）代表未发现，“0”代表实测结果确为零。

如数据已绘成曲线图，可不再列表。

\subsubsection{数学、物理和化学式}

正文中的公式、算式或方程式等应编排序号，序号标注于该式所在行(当有续行时，应
标注于最后一行)的最右边。

较长的式，另行居中横排。如式必须转行时，只能在$+$，$-$，$\times$，$\div$，$<$，
$>$处转行。上下式尽可能在等号“$=$”处对齐。

小数点用“$.$”表示。大于$999$的整数和多于三位数的小数，一律用半个阿拉伯数字符的小
间隔分开，不用千位撇。对于纯小数应将$0$列于小数点之前。

示例：应该写成$94\ 652.023\ 567$和$0.314\ 325$, 不应写成$94,652.023,567$和
$.314,325$。

应注意区别各种字符，如：拉丁文、希腊文、俄文、德文花体、草体；罗马数字和阿拉伯数
字；字符的正斜体、黑白体、大小写、上下脚标（特别是多层次，如“三踏步”）、上下偏差
等。

\subsubsection{计量单位}

报告、论文必须采用国务院发布的《中华人民共和国法定计量单位》，并遵照《中华人
民共和国法定计量单位使用方法》执行。使用各种量、单位和符号，必须遵循附录B所
列国家标准的规定执行。单位名称和符号的书写方式一律采用国际通用符号。

\subsubsection{符号和缩略词}

符号和缩略词应遵照国家标准的有关规定执行。如无标准可循，可采纳本学科或本专业
的权威性机构或学术团体所公布的规定；也可以采用全国自然科学名词审定委员会编印
的各学科词汇的用词。如不得不引用某些不是公知公用的、且又不易为同行读者所理解
的、或系作者自定的符号、记号、缩略词、首字母缩写字等时，均应在第一次出现时一
一加以说明，给以明确的定义。

\subsection{结论}

报告、论文的结论是最终的、总体的结论，不是正文中各段的小结的简单重复。结论应
该准确、完整、明确、精炼。在结论中要清楚地阐明论文中有那些自己完成的成果，特
别是创新性成果；

如果不可能导出应有的结论，也可以没有结论而进行必要的讨论。可以在结论或讨论中
提出建议、研究设想、仪器设备改进意见、尚待解决的问题等。

\subsection{致谢}

可以在正文后对下列方面致谢：

\begin{itemize}
\item 国家科学基金、资助研究工作的奖学金基金、合作单位、资助或支持的企业、组织或个
人；
\item 协助完成研究工作和提供便利条件的组织或个人；
\item 在研究工作中提出建议和提供帮助的人；
\item 给予转载和引用权的资料、图片、文献、研究思想和设想的所有者；
\item 其他应感谢的组织或个人。
\end{itemize}

\subsection{参考文献表}

\subsubsection{专著著录格式}

主要责任者，其他责任者，书名，版本，出版地：出版者，出版年

例：1. 刘少奇，论共产党员的修养，修订2版，北京：人民出版社，1962

\subsubsection{连续出版物中析出的文献著录格式}

析出文献责任者，析出文献其他责任者，析出题名，原文献题名，版本：文献中的位置。

例：2. 李四光，地壳构造与地壳运动，中国科学，1973 (4)：400－429

参考文献采用顺序编码制，按论文正文所引用文献出现的先后顺序连续编码。

\section{附录}

附录是作为报告、论文主体的补充项目，并不是必需的。

下列内容可以作为附录编于报告、论文后，也可以另编成册；

\begin{enumerate}
\item 为了整篇论文材料的完整，但编入正文又有损于编排的条理和逻辑性，这一材料
包括比正文更为详尽的信息、研究方法和技术更深入的叙述，建议可以阅读的参考文献
题录，对了解正文内容有用的补充信息等；
\item 由于篇幅过大或取材于复制品而不便于编入正文的材料；
\item 不便于编入正文的罕见珍贵资料；
\item 对一般读者并非必要阅读，但对本专业同行有参考价值的资料；
\item 某些重要的原始数据、数学推导、计算程序、框图、结构图、注释、统计表、计
算机打印输出件等。
\end{enumerate}

附录与正文连续编页码。

每一附录均另页起。

\section{结尾部分 (必要时)}

为了将论文迅速存储入电子计算机，可以提供有关的输入数据。可以编排分类索引、著者索
引、关键词索引等。

% 参考文献。应放在\backmatter之前。
% 推荐使用BibTeX，若不使用BibTeX时注释掉下面一句。
\nocite{*}
\bibliography{sample}
% 不使用 BibTeX
%\begin{thebibliography}{2}
%
%\bibitem{deng:01a}
%{邓建松,彭冉冉,陈长松}.
%\newblock {\em \LaTeXe{}科技排版指南}.
%\newblock 科学出版社,书号:7-03-009239-2/TP.1516, 北京, 2001.
%
%\bibitem{wang:00a}
%王磊.
%\newblock {\em \LaTeXe{}插图指南}.
%\newblock 2000.
%\end{thebibliography}

% 附录，必须放在参考文献后，backmatter前

%%%%%%%%%%%%%%%%%%%%%%%%%%%%%%%%%%%%%%%%%%%%%%%%%%%%%%%%%%%%%%%%%%%%%%%%%%%%%%%
% 书籍附件
\backmatter
%%%%%%%%%%%%%%%%%%%%%%%%%%%%%%%%%%%%%%%%%%%%%%%%%%%%%%%%%%%%%%%%%%%%%%%%%%%%%%%
% 作者简历与科研成果页，应放在backmatter之后
\begin{resume}
% 论文作者身份简介，一句话即可。
\begin{authorinfo}
\noindent 王东宇，男，汉族，1996年3月出生，海南省临高人。
\end{authorinfo}
% 论文作者教育经历列表，按日期从近到远排列，不包括将要申请的学位。
\begin{education}
\item[2018年9月 --- 2021年6月] 南京大学计算机科学与技术系 \hfill 硕士
\item[2014年9月 --- 2018年6月] 南京大学计算机科学与技术系 \hfill 本科
\end{education}
% 论文作者在攻读学位期间所发表的文章的列表，按发表日期从近到远排列。
\begin{publications}
\item Xiaobao Wei, Jinnan Chen, ``Voting-on-Grid Clustering for Secure
  Localization in Wireless Sensor Networks,'' in \textsl{Proc. IEEE International
    Conference on Communications (ICC) 2010}, May. 2010.
\item Xiaobao Wei, Shiba Mao, Jinnan Chen, ``Protecting Source Location Privacy
  in Wireless Sensor Networks with Data Aggregation,'' in \textsl{Proc. 6th
    International Conference on Ubiquitous Intelligence and Computing (UIC)
    2009}, Oct. 2009.
\end{publications}
% 论文作者在攻读学位期间参与的科研课题的列表，按照日期从近到远排列。
\begin{projects}
\item 国家自然科学基金面上项目``无线传感器网络在知识获取过程中的若干安全问题研究''
（课题年限~2010年1月 --- 2012年12月），负责位置相关安全问题的研究。
\item 江苏省知识创新工程重要方向项目下属课题``下一代移动通信安全机制研究''
（课题年限~2010年1月 --- 2010年12月），负责LTE/SAE认证相关的安全问题研究。
\end{projects}
\end{resume}

%%%%%%%%%%%%%%%%%%%%%%%%%%%%%%%%%%%%%%%%%%%%%%%%%%%%%%%%%%%%%%%%%%%%%%%%%%%%%%%
% 生成《学位论文出版授权书》页面，应放在最后一页
% \makelicense

%%%%%%%%%%%%%%%%%%%%%%%%%%%%%%%%%%%%%%%%%%%%%%%%%%%%%%%%%%%%%%%%%%%%%%%%%%%%%%%
\end{document}
